\begin{frame}{CAP y PACELC: lo decide lo que está detrás del servicio}
\begin{columns}[T]
\begin{column}{0.34\textwidth}
\small
\begin{itemize}
\item Un servicio REST no guarda nada entre peticiones \cite{fielding2000}: decide la base de datos que tiene detrás.
\item Con partición se elige entre disponibilidad y consistencia; sin partición, entre latencia y consistencia \cite{abadi2012}.
\end{itemize}
\end{column}
\begin{column}{0.62\textwidth}
\footnotesize
\begin{tabularx}{\linewidth}{@{}l X@{}}
\toprule
Sistema & Qué prioriza \\
\midrule
DynamoDB & Lectura eventual por defecto, que cuesta la mitad; con \texttt{ConsistentRead=true} devuelve el dato más reciente \cite{aws_consistencia} \\
Caché HTTP & Entrega la copia guardada mientras siga fresca; con \texttt{ETag} pregunta si cambió \cite{rfc9111} \\
Nuestro PFC & Una sola base PostgreSQL: lee lo último, pero si la base cae el servicio falla (C antes que A) \\
\bottomrule
\end{tabularx}
\end{column}
\end{columns}
\begin{block}{CAP no es elegir dos de tres para siempre}
Las particiones de red son raras: casi siempre se puede tener C y A, y la decisión real es latencia o consistencia \cite{abadi2012}.
\end{block}
\end{frame}
